Between the quotient and the multiplier form there is a hybrid, in which the multiplier acts on the kernel directions only and the elimination on the rest.
  We choose a complement $C$ of $\ker \sprod{2}{2}$ in $\Hn{2}$, write $\pi_N$ for the projection on the kernel along $C$ and $\vertex{5 \cdot 5}^{[C]}$ for the composite vertex~\eqref{p:eq:hs5:exchange-vertex} with $\sprod{2}{2}^{-1}$ replaced by the inverse of $\sprod{2}{2}$ restricted to $C$, extended by zero on the annihilator of $C$.
  Since the values of $\sprod{3}{2}$ no longer lie in the image of $\sprod{2}{2}$ when~\eqref{p:eq:hs5:kernel-condition} fails, this extension is where the complement enters.
  We take $\Sigma_2$ on the whole of $\Hn{2}$, with the multiplier $\Theta_2$ of \cref{p:app:conventions} paired with its kernel component only:
  \begin{equation}
  \label{p:eq:hs6:dressed-vertex-C}
  \begin{aligned}
    \dvertex{6}^{[C]}
      &
      := \vertex{6} - \frac{6}{5} \, \vertex{5 \cdot 5}^{[C]},
    \\
    \caction[\le 6]^{[C]}
      &
      := \caction[\le 5]
        + \frac{\gs^4}{6!} \, \dvertex{6}^{[C]}(\Psi \otimes \Sigma_2, \Psi \otimes \Sigma_2)
    \\ &
    \quad
        - \frac{\gs^4}{6!} \, \dvertex{6}^{[C]}(\Sigma_3 - \Psi \otimes \Sigma_2, \Sigma_3 - \Psi \otimes \Sigma_2)
        + \pair[2]{\Theta_2}{\pi_N(\Sigma_2 - \Psi^2)}.
  \end{aligned}
\end{equation}
  Once $\Theta_2$, $\Sigma_3$ and $\Sigma_2$ are eliminated, this action reproduces $S_{\le 6}$ through order $\gs^4$ for every choice of complement, as an identity in $\Psi$ off-shell.
  The classical value of the multiplier starts at order $\gs^3$, and it is nonzero exactly when the quintic vertex does not vanish on the kernel.
  The vertex $\dvertex{6}^{[C]}$ depends on the complement (\cref{p:sec:all:multiplier} and \cref{sec:mult:hybrid-derivation} discuss how a coupling of the multiplier can carry this dependence instead), but the eliminated action does not, through order $\gs^4$; the remainder of order $\gs^5$ does, until the next reduction.
  We established this in the finite-dimensional models of \cref{p:app:finite-models} for the form with $\Sigma_2$ alone, in which the two sextic terms above are replaced by the single term $\dvertex{6}^{[C]}(\Sigma_2^3)$ with the coefficient of the first one, with a kernel of dimension one and two.

  It holds for the string vertices if the exchange terms are not linear combinations of the elementary vertices and the quartic form is invertible on $C$.
  When the quintic vertex vanishes on the kernel, the inverse on $C$ agrees with the inverse on the quotient on all the values of $\sprod{3}{2}$, the dependence on $C$ disappears, and the quotient form above is recovered.
  We have not worked out the gauge transformation of the hybrid, its Batalin--Vilkovisky structure or its closure at order $\gs^5$.
